\chapter{Conclusion}
\label{chap:conclusion}

The four objectives were completed in simulation.
Objectives O1 and O2 produced a maintainable ROS~2 Jazzy and Gazebo Harmonic
environment with selective robot launch, calibrated rail cells, identifiable
shuttles and switch-dependent routes. One all-robot cold start confirmed access
to the six configured models, and all 96 transport checks passed across the
geometry, topology, coordination, device and identity functions.

Objectives O3 and O4 extended this foundation with typed command and state
interfaces, task planning, language interpretation and paired-camera
perception. Together, these elements formed the evaluated English-to-motion
workflow.

The language interface also produced the intended outcome in all ten evaluation
cases, including clarification, rejection and cancellation.

The complete path from an English request to shuttle motion was exercised in
12 isolated cold-start runs. Every run reached its intended terminal state, and
all 24 recorded step postconditions were satisfied. On the
locked 256-scenario Test set, the configured visual checkpoint achieved 100\%
rail-side accuracy, 99.2537\% loaded-state accuracy and 65.2985\% exact-block
accuracy.

For MFJA, the project provides a reusable platform for practical teaching,
software integration, synthetic-data generation and further planning or
perception experiments. For me, the most important lesson was learning how to
divide an initially broad integration request into smaller parts that I could
implement and test one at a time. I also prepared the configuration, regression
tests and run instructions so that the host team can continue using and
extending the work.

The immediate next step is to repeat the campaign across controlled seeds and
conduct a live-graph fault-injection study before extending the simulated
workflow to broader scenarios.
